% Physical unit 008; accepted Session 02.
\GordanUnitStart{008}{printed page 3}
\GordanTopFolio{3}
so on the spheroid we wish to consider more precisely the arc \gq{$s$} of a curve (which, because $\psi$ is the longitude, we choose to call the extension), and the angle \gq{$T$} (azimuth), which the curve makes with the meridian at the points where it cuts the meridian.

\begin{center}\S\ 9.\end{center}
It remains, then, to refer the quantities $s$ and $T$ to the quantities $l$ and $\psi$. Let us therefore consider on some spheroidal curve two points $A(l,\psi)$ and $B(l+dl,\psi+d\psi)$, so close together that no nearer points can be conceived. Let the meridian of the point $A$ cut the circle of the point $B$ in the point $C(l+dl,\psi)$. Since the triangle $ABC$ formed by these points is infinitesimal, it may be taken as lying in a plane. Its sides $AC$ and $BC$, the one tangent to a meridian and the other to a circle, are inclined at a right angle.

Accordingly,
\[
 {BC\over AC}=\tg BAC;\qquad {AC\over AB}=\cos BAC;\qquad {BC\over AB}=\sin BAC,
\]
and
\[
  AB^2=AC^2+BC^2.
\]
But
\[
\begin{gathered}
 BAC=T;\qquad AB=ds,\\
 BC\;\hbox{(the arc of a circle described with radius \gq{$a\cos l$})}=a\cos l\,d\psi,\\
 AC\;\hbox{(the arc of the ellipse)}=a\sqrt{1-e^2\cos^2 l}\,dl.
\end{gathered}
\]
Putting these values in place of the preceding ones, we obtain
\[
 ds^2=a^2(1-e^2\cos^2l)\,dl^2+a^2\cos^2l\,d\psi^2,
\]
\[
\tg T={\cos l\over\sqrt{1-e^2\cos^2l}}{d\psi\over dl};\qquad
\cos T=dl\sqrt{\frac{1-e^2\cos^2l}{(1-e^2\cos^2l)dl^2+\cos^2l\,d\psi^2}},
\]
\[
\sin T=d\psi\sqrt{\frac{\cos^2l}{(1-e^2\cos^2l)dl^2+\cos^2l\,d\psi^2}}.
\]

\begin{center}\S\ 10.\end{center}
For the earth, from the calculations proposed by the illustrious Bessel, the following ratios are relevant:
\[
\begin{array}{rlrl}
 a&=3272077,1399\ \hbox{toises}; & b&=3261139,3284\ \hbox{toises},\\
 \log a&=6,5148235337; & \log b&=6,5133693539,\\
 \log {b\over a}&=9,9985458202; & \log {a\over b}&=0,0014541798,\\
 \varepsilon={a-b\over a+b}&={1-\sqrt{1-e^2}\over 1+\sqrt{1-e^2}}=0,001674184767,\\
 \log e&=8,9122052075; & \log\varepsilon&=7,2238033861.
\end{array}
\]
